\section{The CTL indistinguishability induction}
\label{sec:annex-p09-ctl-witnesses}

\begin{annexitem}{annex:p09-bk-ctl-witness-depth}{thm:p09-bk-ctl-witness-depth}{Bounded CTL depth in the Baier--Katoen witnesses}
\begin{proof}
All systems here are finite and total. The graph below the top pair
of \(TS_n\) or \(TS'_n\), for \(n>0\), is the same closed copy of
\(TS'_{n-1}\); no edge leaves that copy. Consequently the truth of
a CTL state formula at one of its states is unaffected by the larger
system in which it is embedded. This follows directly by induction
on state formulas, since the available paths from the state are the
same. We may therefore compare the states of different levels without
renaming the lower copies.

It suffices to use the existential CTL basis consisting of atoms,
Boolean operations, \(EX\), \(EG\), and \(E[\alpha U\beta]\).
The usual eliminations preserve temporal depth. Explicitly,
\[
\begin{aligned}
 AX\alpha&\equiv\neg EX\neg\alpha,
 &AF\alpha&\equiv\neg EG\neg\alpha,\\
 EF\alpha&\equiv E[\mathsf{true}\,U\alpha],
 &AG\alpha&\equiv\neg E[\mathsf{true}\,U\neg\alpha],\\
 A[\alpha U\beta]
 &\equiv\neg E[\neg\beta\,U(\neg\alpha\land\neg\beta)]
          \land\neg EG\neg\beta.
\end{aligned}
\]
For the last identity, failure of universal strong until means either
that a path never reaches \(\beta\), or that before its first
\(\beta\)-state it reaches a state satisfying neither \(\alpha\)
nor \(\beta\). Conversely either alternative falsifies universal
until. Totality allows any finite witness path to be continued
infinitely. The other identities follow immediately from the path
clauses. In these displays a path quantifier and its temporal operator
count as one level, as in the statement.

We prove simultaneously, by structural induction on a formula,
its eventual constant truth value on each of the two primed state
families and its equality at the primed and unprimed top states above
the asserted depth. Atoms depend only on the two labels: every
\(s\)-state satisfies \(\neg a\), and every \(t\)-state satisfies
\(a\). The Boolean steps follow by applying the induction hypothesis
to each operand and taking the maximum of their depths.

\emph{Next.} Let \(\Phi=EX\alpha\) and \(k=d(\Phi)\).
For \(m\geq k\), the only successor of \(s'_m\) is \(t'_m\),
and the successors of \(t'_m\) are \(t'_m,s'_{m-1}\).
The truth values of \(\alpha\) on these states are constant by
induction, because \(m-1\geq k-1=d(\alpha)\). Hence both primed
truth values stabilize. The unprimed \(s_m\) again has only its
\(t_m\) successor. The unprimed \(t_m\) has the additional successor
\(s_m\), but \(\alpha\) has the same truth value at \(s_m\) as
at \(s'_m\) and \(s'_{m-1}\). The extra successor therefore changes
neither existential-next truth value.

\emph{Existential always.} At any primed or unprimed \(t\)-state,
\(EG\alpha\) holds exactly when \(\alpha\) holds there: the
self-loop supplies a witnessing path, and the current state must
satisfy \(\alpha\). At the associated \(s\)-state it holds exactly
when \(\alpha\) holds both at that state and at its unique
\(t\)-successor; afterwards the same self-loop suffices. The
induction hypothesis proves both required comparisons.

\emph{Existential until: stabilization on the primed systems.}
Let \(\Phi=E[\alpha U\beta]\) and
\(k=1+\max\{d(\alpha),d(\beta)\}\). For each \(i\), let
\(x_i,y_i\in\{0,1\}\) be the truth values of \(\Phi\) at
\(s'_i,t'_i\). Write \(a_s(i),a_t(i)\) for the truth values of
\(\alpha\) at those states and \(b_s(i),b_t(i)\) for those of
\(\beta\). Use \(\land,\lor\) also for Boolean operations on
these truth values. A witness to strong until is a finite path ending
at a \(\beta\)-state, all preceding states satisfying \(\alpha\).
Repeated self-loops before the final state may be deleted from such
a witness. Thus at the base level
\[
 y_0=b_t(0),\qquad
 x_0=b_s(0)\lor\bigl(a_s(0)\land b_t(0)\bigr),
\]
and for every \(i>0\),
\[
\begin{aligned}
 y_i&=b_t(i)\lor\bigl(a_t(i)\land x_{i-1}\bigr),\\
 x_i&=b_s(i)\lor\bigl(a_s(i)\land y_i\bigr).
\end{aligned}
\]
The equations include a witness of length zero. They use every possible
non-loop successor, so they give equivalences, not merely sufficient
conditions.

By induction, all four coefficient values are constant for
\(i\geq k-1\); denote them by \(a_s,a_t,b_s,b_t\). Put
\[
 B=b_s\lor(a_s\land b_t),\qquad A=a_s\land a_t,
 \qquad g(z)=B\lor(A\land z).
\]
The elementary Boolean identity \(g(g(z))=g(z)\) is obtained by
distributing \(A\land\) over the disjunction and absorbing
\(A\land B\) into \(B\). If \(k\geq2\), the displayed recurrence
gives \(x_i=g(x_{i-1})\) for every \(i\geq k-1\), hence
\(x_i=x_{i+1}\) throughout that range. If \(k=1\), the base
equation says \(x_0=B\), and \(g(B)=B\) gives the same conclusion
from level zero. In either case \(x_i\) is constant for
\(i\geq k-1\), and the equation for \(y_i\) makes \(y_i\)
constant for \(i\geq k\). This is the required primed stabilization.

\emph{Existential until: the extra return edge.} Fix \(n\geq k\).
Induction gives the same operand values \(a_s,a_t,b_s,b_t\) at
\(s_n,t_n\) as at their primed counterparts. Let \(L=x_{n-1}\)
be the until truth value at the common lower state \(s'_{n-1}\).
For the unprimed top pair, let \(X,Y\) be its until truth values.
Deleting self-loops and repeated circuits through the top pair from
a finite witness gives
\[
 Y=b_t\lor(a_t\land b_s)\lor(a_t\land L),
 \qquad X=b_s\lor(a_s\land Y).
\]
Indeed, a witness from \(t_n\) can finish there, can next finish
at \(s_n\), or can descend into the lower copy; any return from
\(s_n\) to \(t_n\) before the endpoint can be erased. Conversely
the three disjuncts respectively supply those three finite witnesses.
For the second disjunct a single transition suffices, so no condition
\(a_s\) is needed.

The induction hypothesis for \(\beta\) makes its value at
\(s'_{n-1}\) equal to \(b_s\), since \(n-1\geq k-1\).
Consequently \(b_s=1\) implies \(L=1\) by the length-zero until
witness. The middle disjunct in \(Y\) is absorbed into the last;
therefore
\[
 Y=b_t\lor(a_t\land L)=y_n,\qquad X=x_n.
\]
This proves the comparison at both top states and completes the
simultaneous induction.

Finally, \((s_nt_n)^\omega\) is a path of \(TS_n\) with infinitely
many \(\neg a\)-states. In \(TS'_n\), any move from a \(t'_i\)
to a different pair strictly decreases the nonnegative level. There
can be only finitely many such moves, so every path eventually stays
at a self-looping \(a\)-state. Thus \(TS_n\) fails the universal
\(\Diamond\Box a\) property and \(TS'_n\) satisfies it. If a CTL
formula \(\Phi\) defined that property, choosing
\(n\geq d(\Phi)\) would contradict the equality at the two initial
states just proved. Every occurrence of a temporal operator contributes
at least one symbol to formula length, so \(d(\Phi)\leq|\Phi|\)
also recovers the source's length-bounded claim.
\end{proof}
\end{annexitem}
